\chapter{Bài toán kiểm tra màn hình điện thoại}
\label{chap:c1}

\section{Mô tả bài toán}

Kiểm tra và phân loại lỗi bề mặt màn hình điện thoại là công đoạn then chốt trong quy trình kiểm soát chất lượng sản phẩm tại các dây chuyền lắp ráp công nghiệp \cite{park2023design}. Cấu trúc phức hợp gồm nhiều lớp kính, tấm nền hiển thị và cảm ứng khiến các khuyết tật xuất hiện với đặc tính quang học đa dạng như vết trầy xước, hạt bụi, điểm chết pixel, vết nứt v.v… Hệ thống kiểm tra quang học tự động hỗ trợ phát hiện kịp thời các khuyết tật phát sinh, góp phần giảm chi phí nhân công và duy trì hiệu suất vận hành toàn hệ thống.

Các phương pháp kiểm tra lỗi truyền thống chủ yếu dựa trên kỹ thuật học máy có giám sát và các mạng nơ-ron tích chập \cite{nguyen2017novel, li2020deep, yolo2023mobile}. Kỹ thuật học máy có giám sát đòi hỏi tập dữ liệu gán nhãn thủ công quy mô lớn đối với từng loại lỗi cố định. Trong thực tế sản xuất, các hình thái lỗi mới phát sinh liên tục do sự thay đổi của linh kiện đầu vào hoặc sự cố kỹ thuật của thiết bị. Do đó, các mô hình học giám sát có cấu trúc cố định thường gặp khó khăn trong việc nhận diện các dạng khuyết tật chưa từng xuất hiện trong tập dữ liệu huấn luyện.

Xu hướng ứng dụng các mô hình học không giám sát trong phát hiện bất thường giúp giảm sự phụ thuộc vào dữ liệu lỗi được gán nhãn thủ công \cite{liu2024effnet, roth2022smartphone}. Các phương pháp này học tập các đặc trưng không gian từ tập sản phẩm đạt chuẩn để xác định ngưỡng sai biệt cho các mẫu khuyết tật. Tuy nhiên, phần lớn các tiếp cận không giám sát hiện hành chưa tích hợp cơ chế phân loại tự động và tổ chức dữ liệu cho các loại lỗi mới xuất hiện theo thời gian thực.

Nghiên cứu này tập trung giải quyết bài toán phát hiện bất thường và phân loại các dạng lỗi bề mặt màn hình điện thoại trong môi trường dòng dữ liệu động. Hệ thống đặt ra ba yêu cầu kỹ thuật bao gồm (i) phát hiện chính xác các vùng ảnh chứa lỗi dựa trên phân phối của mẫu sản phẩm chuẩn, (ii) tự động phân cụm các dạng lỗi mới phát sinh mà không làm gián đoạn tiến trình vận hành của hệ thống, và (iii) duy trì khả năng học liên tục để cập nhật tri thức mới. Dữ liệu hình ảnh thu thập từ thực tế gặp hiện tượng mất cân bằng nghiêm trọng giữa số lượng mẫu đạt chuẩn và mẫu lỗi. Nghiên cứu đề xuất giải pháp tích hợp giữa kỹ thuật phát hiện bất thường ngoại tuyến và thuật toán học tăng cường trực tuyến nhằm phân loại các cụm lỗi theo thời gian thực.

\begin{figure}[htbp]
\centering
\begin{tikzpicture}[node distance=3.2cm, auto, >=stealth,
    block/.style={rectangle, draw=blue!80!black, fill=blue!5!white, text width=3.2cm, text centered, minimum height=1.2cm, font=\normalsize, rounded corners=2pt},
    line/.style={draw=blue!80!black, ->, thick}]
    \node [block] (input) {Ảnh màn hình điện thoại};
    \node [block, right of=input, xshift=2.5cm] (phase1) {Phát hiện bất thường ngoại tuyến};
    \node [block, right of=phase1, xshift=2.5cm] (phase2) {Phân cụm lỗi trực tuyến};
    \node [block, below of=phase2, yshift=-0.5cm] (output) {Phân loại dạng lỗi};

    \path [line] (input) -- (phase1);
    \path [line] (phase1) -- node[above, font=\normalsize, text=black] {Ảnh lỗi} (phase2);
    \path [line] (phase2) -- (output);
\end{tikzpicture}
\caption{Sơ đồ tổng quan quy trình kiểm tra và phân loại lỗi màn hình điện thoại.}
\label{fig:quy_trinh_tong_quan}
\end{figure}

Sơ đồ tổng quan quy trình kiểm tra và phân loại lỗi màn hình điện thoại được mô tả cụ thể trong Hình \ref{fig:quy_trinh_tong_quan}. Luồng xử lý bắt đầu từ việc tiếp nhận ảnh bề mặt màn hình điện thoại, sau đó dữ liệu được đưa qua khối phát hiện bất thường ngoại tuyến để lọc ra các mẫu sản phẩm đạt chuẩn và giữ lại các ảnh chứa vị trí lỗi. Các ảnh chứa lỗi tiếp tục được chuyển sang khối phân cụm lỗi trực tuyến nhằm tự động gom nhóm và xác định chủng loại khuyết tật mới.

Tiêu chuẩn đánh giá một mẫu ảnh thuộc nhóm bất thường được thiết lập dựa trên khoảng cách biểu diễn không gian ẩn so với ngưỡng phân phối mẫu chuẩn, được xác định thông qua Công thức \ref{eq:threshold}:
\begin{equation}
d(z, \mathcal{W}) > 3\sigma
\label{eq:threshold}
\end{equation}
Trong Công thức \ref{eq:threshold}, ký hiệu $z$ biểu thị véc-tơ biểu diễn ẩn thu được từ không gian đặc trưng của mẫu ảnh đầu vào, ký hiệu $\mathcal{W}$ đại diện cho tập hợp các véc-tơ trọng số đại diện cho phân phối mẫu sản phẩm đạt chuẩn trong mô hình huấn luyện, ký hiệu $d(z, \mathcal{W})$ thể hiện khoảng cách Euclid nhỏ nhất từ véc-tơ ẩn $z$ tới tập hợp trọng số $\mathcal{W}$, và ký hiệu $\sigma$ đại diện cho độ lệch chuẩn của sai số tái tạo trên tập dữ liệu mẫu đạt chất lượng. Tất cả các mẫu ảnh có giá trị khoảng cách vượt quá ngưỡng phân phối này sẽ được xác định là bất thường và được chuyển sang khối phân cụm trực tuyến.

Giải pháp đề xuất sử dụng mạng Growing Neural Gas cải tiến kết hợp với cơ chế tích lũy bằng chứng nhằm tối ưu hóa tiến trình phân loại trực tuyến \cite{fritzke1995growing}. Phương pháp học động dựa trên mạng Growing Neural Gas có khả năng tự động thay đổi cấu trúc đồ thị và cập nhật tri thức mới mà không gặp phải hiện tượng quên tri thức cũ. Nghiên cứu kết hợp kỹ thuật tính toán bản đồ entropy và mô hình mã hóa Autoencoder trong pha phát hiện bất thường ngoại tuyến để trích xuất đặc trưng không gian và giảm chiều dữ liệu. Trong pha phân loại trực tuyến, thuật toán cải tiến Evidence-Accumulating Growing Neural Gas kiểm soát mật độ nút và tự động gộp các cụm lỗi gần nhau, giúp nâng cao độ chính xác phân loại trong hệ thống kiểm tra công nghiệp.
\section{Ý tưởng và hướng tiếp cận của bài toán}

Ý tưởng chủ đạo của luận văn dựa trên việc chia tách quá trình kiểm tra ngoại quan thành hai pha độc lập bao gồm pha phát hiện bất thường ngoại tuyến và pha phân cụm lỗi trực tuyến. Hướng tiếp cận này giải quyết hai thách thức lớn trong kiểm tra chất lượng màn hình điện thoại bao gồm sự mất cân bằng nghiêm trọng giữa số lượng mẫu đạt chuẩn và mẫu lỗi, cùng khả năng tự thích ứng với các dạng khuyết tật mới phát sinh trong dây chuyền sản xuất \cite{goodfellow2016deep}. Trong pha thứ nhất, hệ thống tập trung xây dựng mô hình biểu diễn không gian cho các mẫu sản phẩm đạt chất lượng nhằm nhanh chóng sàng lọc và loại bỏ phần lớn dữ liệu bình thường. Trong pha thứ hai, hệ thống áp dụng cơ chế học động trực tuyến để tự động gom nhóm và phân loại các mẫu bị lỗi mà không cần gán nhãn trước \cite{fritzke1995growing}.

Pha phát hiện bất thường ngoại tuyến kết hợp kỹ thuật trích xuất thông tin kết cấu không gian và mô hình học biểu diễn không giám sát. Dữ liệu ảnh đầu vào ban đầu gồm ba kênh màu chuẩn được bổ sung một kênh bản đồ entropy cục bộ \cite{shannon1948mathematical}. Việc ghép nối kênh bản đồ này tạo thành cấu trúc ảnh bốn kênh nhằm tăng cường khả năng nhận diện các vùng có sự thay đổi đột ngột về độ sáng hoặc bề mặt. Tiếp theo, mô hình mạng mã hóa tự động thực hiện nén dữ liệu ảnh bốn kênh thành các véc-tơ biểu diễn ẩn mang thông tin cốt lõi của sản phẩm đạt chuẩn. Mô hình Mạng khí Nơ-ron tăng trưởng ngoại tuyến được huấn luyện trên tập các véc-tơ biểu diễn ẩn này để thiết lập bản đồ ma trận trọng số đại diện cho phân phối dữ liệu bình thường. Một mẫu ảnh được xác định là bất thường khi véc-tơ biểu diễn ẩn tương ứng có khoảng cách không gian vượt quá ngưỡng biến động quy định so với các nút trọng số thuộc mô hình Mạng khí Nơ-ron Tăng trưởng.

Pha phân cụm lỗi trực tuyến xử lý các mẫu ảnh chứa khuyết tật đã được sàng lọc từ pha thứ nhất thông qua mạng khí nơ-ron tăng trưởng. Mô hình ResNet được sử dụng để trích xuất các véc-tơ đặc trưng tập trung vào hình thái khuyết tật từ cấu trúc ảnh bốn kênh \cite{he2016deep}. Các véc-tơ đặc trưng này được truyền trực tiếp vào thuật toán Mạng khí Nơ-ron tăng trưởng Tích lũy Bằng chứng để cập nhật cấu trúc mạng theo thời gian thực. Thuật toán cải tiến tích hợp bốn cơ chế cốt lõi bao gồm (i) điều chỉnh tốc độ học và sai số tích lũy theo tuổi của nút, (ii) chèn nút muộn dựa trên vùng đệm bằng chứng và huấn luyện cục bộ, (iii) kiểm soát sự phát triển hình thái đồ thị dựa trên mật độ nút, và (iv) kiểm tra sự chồng lấp hộp giới hạn giữa các cụm để thực hiện gộp vết. Cuối cùng, thuật toán phân cụm thành phần liên thông trên đồ thị được áp dụng để xác định chính xác các nhóm lỗi riêng biệt mà không cần xác định trước số lượng cụm.

\begin{figure}[htbp]
\centering
\begin{tikzpicture}[node distance=1.5cm, auto, >=stealth,
    block/.style={rectangle, draw, fill=white, text width=3.3cm, text centered, minimum height=1.2cm, font=\normalsize},
    process/.style={rectangle, draw, fill=white, text width=3.3cm, text centered, minimum height=1.2cm, font=\normalsize},
    subblock/.style={rectangle, draw, fill=white, text width=3.3cm, text centered, minimum height=1.0cm, font=\normalsize},
    line/.style={draw, ->, thick}]
    
    \node [block] (input) {Ảnh màn hình RGB};
    \node [process, right of=input, xshift=4.0cm] (entropy) {Trích xuất Entropy và tạo ảnh 4 kênh};
    \node [process, right of=entropy, xshift=4.0cm] (ae) {Autoencoder trích xuất véc-tơ ẩn $z$};
    \node [subblock, below of=ae, yshift=-2cm] (gng_off) {Offline GNG phát hiện bất thường};
    
    \node [process, below of=entropy, yshift=-2cm] (resnet) {ResNet trích xuất đặc trưng hình thái};
    \node [subblock, left of=resnet, xshift=-4.0cm] (eagng) {Online EAGNG phân cụm lỗi thích ứng};
    \node [block, below of=eagng, yshift=-2cm] (cluster) {Phân loại lỗi bằng thành phần liên thông};

    \path [line] (input) -- (entropy);
    \path [line] (entropy) -- (ae);
    \path [line] (ae) -- (gng_off);
    \path [line] (gng_off) -- node[above, font=\normalsize] {Ảnh lỗi} (resnet);
    \path [line] (resnet) -- (eagng);
    \path [line] (eagng) -- (cluster);
\end{tikzpicture}
\caption{Sơ đồ kiến trúc hướng tiếp cận hai pha trong kiểm tra và phân loại lỗi màn hình điện thoại.}
\label{fig:huong_tiep_can}
\end{figure}

Sơ đồ kiến trúc xử lý chi tiết được mô tả cụ thể trong Hình \ref{fig:huong_tiep_can}. Luồng dữ liệu vận hành khép kín từ khâu tiếp nhận ảnh gốc ba kênh màu, tiến hành tính toán bản đồ độ hỗn loạn không gian để chuyển đổi thành cấu trúc bốn kênh, sau đó đưa qua mạng mã hóa tự động để trích xuất không gian ẩn. Khối Mạng khí Nơ-ron tăng trưởng ngoại tuyến đóng vai trò bộ lọc phân phối để loại bỏ các mẫu sản phẩm đạt chất lượng. Các mẫu ảnh vượt ngưỡng phân phối được chuyển tiếp sang khối ResNet để trích xuất đặc trưng hình thái khuyết tật, làm đầu vào cho mô hình Mạng khí Nơ-ron tăng trưởng Tích lũy Bằng chứng trực tuyến thực hiện tích lũy và điều chỉnh cấu trúc đồ thị trước khi phân loại cụ thể.

Quá trình trích xuất kênh bản đồ độ hỗn loạn thông tin được thực hiện dựa trên việc đo lường mức độ biến động tần suất xuất hiện của các giá trị mức xám trong vùng lân cận cục bộ của từng điểm ảnh \cite{shannon1948mathematical}. Chỉ số này phản ánh mức độ không đồng nhất về bề mặt, giúp làm nổi bật các đốm tròn mờ, điểm chết hoặc bọt khí trên tấm nền hiển thị mà các kênh màu thông thường khó phát hiện. Sau đó, mạng mã hóa tự động nén thông tin từ ảnh bốn kênh thành không gian ẩn dung lượng thấp thông qua quá trình tối thiểu hóa sai số tái tạo giữa ảnh đầu vào và ảnh đầu ra. Các véc-tơ ẩn thu được chứa đựng các đặc trưng hình học đại diện cho sản phẩm đạt chuẩn, phục vụ cho việc xây dựng ranh giới phân phối của dữ liệu bình thường.

Mô hình Mạng khí Nơ-ron tăng trưởng Tích lũy Bằng chứng quản lý quá trình cập nhật đồ thị thông qua cơ chế suy giảm tốc độ học theo độ tuổi của từng nút. Các nút mới được tạo ra có tốc độ học cao để nhanh chóng di chuyển về vị trí trung tâm của vùng dữ liệu mới phát sinh. Khi độ tuổi của nút tăng lên theo số bước thời gian, tốc độ học sẽ giảm dần theo hàm mũ nhằm giúp vị trí của nút đạt trạng thái cân bằng và duy trì độ ổn định cấu trúc. Kết hợp với việc tích lũy bằng chứng dữ liệu trong vùng đệm trước khi quyết định chèn nút, hướng tiếp cận này đảm bảo hệ thống phản ứng linh hoạt với các dạng lỗi mới mà không bị ảnh hưởng bởi nhiễu tín hiệu ngẫu nhiên v.v…
\section{GNG trong kiểm tra màn hình điện thoại}

\subsection{Tổng quan về GNG}

Thuật toán mạng khí nơ-ron tăng trưởng (Growing Neural Gas - GNG) đại diện cho mô hình học không giám sát có khả năng tự thích ứng cấu trúc đồ thị không gian dựa trên phân phối dữ liệu đầu vào \cite{fritzke1995growing}. Khác với các thuật toán phân cụm cố định số lượng trung tâm, mô hình khởi đầu với một số lượng nút tối thiểu và tự động bổ sung các nút mới trong tiến trình huấn luyện nhằm tối ưu hóa độ bao phủ không gian đặc trưng. Cấu trúc mạng bao gồm tập hợp các nút mang vectơ trọng số và tập hợp các cạnh nối thiết lập mối quan hệ lân cận giữa các nút.

Nguyên lý vận hành của mạng dựa trên việc xác định hai nút có khoảng cách nhỏ nhất so với véc-tơ dữ liệu đầu vào trong không gian đặc trưng. Nút gần nhất được xác định là nút chiến thắng thứ nhất, trong khi nút gần thứ hai được xác định là nút chiến thắng thứ hai. Mô hình thực hiện điều chỉnh véc-tơ trọng số của nút chiến thắng thứ nhất cùng các nút lân cận kết nối trực tiếp về phía dữ liệu mới. Đồng thời, độ tuổi của cạnh nối giữa hai nút chiến thắng được khởi tạo lại về giá trị không, trong khi độ tuổi của các cạnh lân cận khác bị tăng lên một đơn vị. Các cạnh vượt quá ngưỡng tuổi quy định sẽ bị xóa khỏi hệ thống, dẫn đến việc loại bỏ các nút không còn kết nối.

Cơ chế tăng trưởng cấu trúc đồ thị được điều phối thông qua tiến trình tích lũy sai số tại các nút chiến thắng. Sau một số bước thời gian cố định, hệ thống xác định nút có tổng sai số tích lũy lớn nhất để tiến hành chèn nút mới vào khoảng giữa nút đó và nút lân cận có sai số cao nhất. Tiến trình này giúp mạng nâng cao độ phân giải biểu diễn tại các vùng không gian có mật độ dữ liệu phức tạp. Sai số của các nút liên quan sau đó được giảm theo một hệ số quy định để duy trì tính cân bằng cho toàn hệ thống.

\subsection{Ứng dụng và các nhược điểm của GNG}

Mô hình Mmạng khí nơ-ron tăng trưởng được áp dụng rộng rãi trong các bài toán phân nhóm không gian, nhận dạng mẫu và nén dữ liệu phi tuyến \cite{martinetz1993neural, fritzke1995growing}. Trong lĩnh vực xử lý ảnh và kiểm tra ngoại quan công nghiệp, thuật toán cung cấp giải pháp dựng bản đồ đặc trưng tự động mà không đòi hỏi giả định ban đầu về phân phối xác suất của dữ liệu. Khả năng tự do điều chỉnh số lượng nút giúp mô hình biểu diễn hiệu quả các không gian dữ liệu có hình thái phức tạp.

Mặc dù sở hữu khả năng thích ứng cấu trúc, mô hình truyền thống tồn tại bốn hạn chế kỹ thuật cốt lõi khi vận hành trong môi trường công nghiệp thực tế. Thứ nhất, thuật toán nhạy cảm với nhiễu dữ liệu dẫn đến việc chèn nút không cần thiết tại các vùng không gian có mật độ dữ liệu thưa thớt. Thứ hai, việc tích lũy sai số liên tục không tính đến thời gian tồn tại của nút gây ra sự bùng nổ số lượng nút tại vùng dữ liệu có mật độ cao, làm gia tăng chi phí tính toán của hệ thống. Thứ ba, tốc độ học cố định làm giảm tính ổn định của mạng ở giai đoạn sau của quá trình huấn luyện trực tuyến, gây ra hiện tượng dao động trọng số xung quanh điểm cân bằng. Thứ tư, thuật toán thiếu cơ chế đánh giá sự chồng lấp hình học giữa các cụm độc lập, dẫn đến hiện tượng kết nối sai cạnh giữa các nhóm dữ liệu riêng biệt.

\subsection{GNG đối với bài toán kiểm tra hình ảnh lỗi trên màn hình điện thoại}

Đặc thụ của bài toán kiểm tra ngoại quan màn hình điện thoại đòi hỏi mô hình học không giám sát phải có khả năng xử lý dòng dữ liệu biến động liên tục và hạn chế ảnh hưởng của nhiễu tín hiệu quang học. Các dạng khuyết tật bề mặt như đốm tròn mờ, điểm chết hoặc bọt khí có kích thước và mật độ phân bố đa dạng trên tấm nền hiển thị. Khi áp dụng mô hình truyền thống vào bài toán này, sự xuất hiện của nhiễu ảnh công nghiệp khiến cấu trúc đồ thị bị biến dạng, gây ra tình trạng gom cụm sai lệch giữa các dạng lỗi khác nhau.

Để giải quyết các hạn chế nêu trên, luận văn định hướng ứng dụng thuật toán Mạng khí Nơ-ron Tăng trưởng Tích lũy Bằng chứng trong pha phân cụm lỗi trực tuyến. Phương pháp cải tiến tích hợp bốn cơ chế kỹ thuật bao gồm (i) điều chỉnh tốc độ học và hệ số giảm sai số theo độ tuổi của nút, (ii) chèn nút muộn dựa trên vùng đệm bằng chứng kết hợp huấn luyện cục bộ, (iii) kiểm soát sự phát triển hình thái dựa trên mật độ nút không gian, và (iv) kiểm tra sự chồng lấp hộp giới hạn giữa các cụm để gộp vết. Các cải tiến này giúp mô hình duy trì độ ổn định cấu trúc đồ thị và nâng cao độ chính xác phân loại khi xử lý dòng dữ liệu hình ảnh lỗi.

Tiến trình xử lý dữ liệu và phát triển cấu trúc mô hình cải tiến được mô tả cụ thể trong Hình \ref{fig:eagng_process}. Khối tiếp nhận véc-tơ đặc trưng đưa dữ liệu đầu vào vào vùng đệm bằng chứng để tích lũy thông tin thống kê trước khi quyết định thay đổi cấu trúc đồ thị. Tiến trình kiểm tra mật độ không gian và kiểm tra chồng lấp hộp giới hạn đảm bảo các nút mới chỉ được tạo ra tại vùng dữ liệu có đủ mật độ, đồng thời ngăn chặn sự hợp nhất sai lệch giữa các cụm lỗi độc lập.

\begin{figure}[htbp]
\centering
\begin{tikzpicture}[node distance=1.6cm, auto, >=stealth,
    block/.style={rectangle, draw, fill=white, text width=3.4cm, text centered, minimum height=1.1cm, font=\normalsize},
    process/.style={rectangle, draw, fill=white, text width=3.4cm, text centered, minimum height=1.1cm, font=\normalsize},
    decision/.style={rectangle, draw, fill=white, text width=3.4cm, text centered, minimum height=1.1cm, font=\normalsize},
    line/.style={draw, ->, thick}]
    
    \node [block] (input) {Véc-tơ đặc trưng lỗi từ ResNet};
    \node [process, right of=input, xshift=3.2cm] (buffer) {Vùng đệm tích lũy bằng chứng};
    \node [decision, right of=buffer, xshift=3.2cm] (density) {Kiểm tra mật độ không gian};
    \node [process, below of=density, yshift=-0.8cm] (overlap) {Kiểm tra chồng lấp hộp giới hạn};
    \node [block, left of=overlap, xshift=-3.2cm] (eagng_node) {Cập nhật nút theo độ tuổi};
    \node [block, left of=eagng_node, xshift=-3.2cm] (clustering) {Phân cụm thành phần liên thông};

    \path [line] (input) -- (buffer);
    \path [line] (buffer) -- (density);
    \path [line] (density) -- (overlap);
    \path [line] (overlap) -- (eagng_node);
    \path [line] (eagng_node) -- (clustering);
\end{tikzpicture}
\caption{Tiến trình xử lý dữ liệu và phát triển cấu trúc mô hình EAGNG trực tuyến.}
\label{fig:eagng_process}
\end{figure}

Sự kết hợp giữa cơ chế tích lũy bằng chứng và kiểm soát mật độ giúp giảm thiểu việc khởi tạo các nút dư thừa do nhiễu môi trường sản xuất. Đồng thời, việc gộp các vết lỗi dựa trên hộp giới hạn hình học cho phép hệ thống nhận diện chính xác hình thái của từng chủng loại khuyết tật trên bề mặt màn hình điện thoại.
\section{Động cơ nghiên cứu}

Mục tiêu tổng quát của luận văn là nghiên cứu, thiết kế và phát triển phương pháp học động hai pha kết hợp giữa mô hình học biểu diễn không gian ngoại tuyến và thuật toán học trực tuyến thích ứng nhằm nâng cao độ chính xác và độ ổn định trong hệ thống kiểm tra chất lượng màn hình điện thoại. Luận văn hướng tới việc giải quyết triệt me hai thách thức lớn bao gồm sự mất cân bằng giữa số lượng mẫu đạt chuẩn và mẫu lỗi, cùng khả năng tự động phân loại các dạng khuyết tật mới phát sinh trên dây chuyền sản xuất. Để đạt được mục tiêu tổng quát, nghiên cứu xác định bốn mục tiêu cụ thể bao gồm (i) xây dựng cơ chế trích xuất đặc trưng đa kênh kết hợp bản đồ độ hỗn loạn không gian và mạng mã hóa tự động, (ii) thiết lập bộ lọc phát hiện bất thường ngoại tuyến dựa trên Mạng khí Nơ-ron Tăng trưởng, (iii) phát triển thuật toán Mạng khí Nơ-ron tăng trưởng Tích lũy Bằng chứng trực tuyến để phân cụm các dạng lỗi mới, và (iv) đánh giá hiệu năng hệ thống trên tập dữ liệu ảnh kiểm tra ngoại quan màn hình điện thoại thực tế v.v…

Nội dung nghiên cứu của luận văn tập trung giải quyết các bài toán kỹ thuật từ khâu xử lý dữ liệu đầu vào đến khâu tổ chức phân cụm trực tuyến. Nghiên cứu triển khai việc trích xuất thông tin độ hỗn loạn không gian để tạo lập cấu trúc ảnh bốn kênh, giúp làm nổi bật các đốm mờ và bọt khí bề mặt. Tiếp theo, luận văn nghiên cứu mô hình nén dữ liệu bằng mạng mã hóa tự động kết hợp thuật toán Mạng khí Nơ-ron tăng trưởng ngoại tuyến nhằm xác định ranh giới phân phối của sản phẩm đạt chất lượng. Đối với pha phân cụm lỗi trực tuyến, nội dung nghiên cứu tập trung đề xuất thuật toán Mạng khí Nơ-ron tăng trưởng Tích lũy Bằng chứng với bốn cải tiến cốt lõi bao gồm cơ chế giảm tốc độ học theo độ tuổi của nút, chèn nút muộn dựa trên vùng đệm tích lũy bằng chứng, kiểm soát sự phát triển đồ thị dựa trên mật độ không gian và gộp vết dựa trên sự chồng lấp hộp giới hạn \cite{fritzke1995growing}. Cuối cùng, luận văn tiến hành thử nghiệm, so sánh đối chứng hiệu năng của phương pháp đề xuất với các thuật toán học không giám sát kinh điển \cite{martinetz1993neural}.

\begin{figure}[htbp]
\centering
\begin{tikzpicture}[node distance=1.6cm, auto, >=stealth,
    block/.style={rectangle, draw, fill=white, text width=3.4cm, text centered, minimum height=1.1cm, font=\normalsize},
    process/.style={rectangle, draw, fill=white, text width=3.4cm, text centered, minimum height=1.1cm, font=\normalsize},
    line/.style={draw, ->, thick}]
    
    \node [block] (obj) {Mục tiêu nghiên cứu: Kiểm tra lỗi tự động};
    \node [process, right of=obj, xshift=3.2cm] (content) {Nội dung: Phát triển thuật toán EAGNG};
    \node [block, right of=content, xshift=3.2cm] (result) {Kết quả: Phân loại lỗi trực tuyến chính xác};

    \path [line] (obj) -- (content);
    \path [line] (content) -- (result);
\end{tikzpicture}
\caption{Mối quan hệ giữa mục tiêu, nội dung nghiên cứu và kết quả dự kiến của luận văn.}
\label{fig:dong_co_nghien_cuu}
\end{figure}

Mối liên hệ logic giữa mục tiêu, nội dung nghiên cứu và kết quả dự kiến được thể hiện cụ thể trong Hình \ref{fig:dong_co_nghien_cuu}. Tiến trình nghiên cứu bắt đầu từ việc xác định mục tiêu kiểm tra lỗi tự động, từ đó triển khai nội dung phát triển thuật toán cải tiến nhằm đạt được kết quả phân loại lỗi trực tuyến chính xác trong môi trường sản xuất công nghiệp.

Kết quả dự kiến của luận văn đóng góp cả về mặt lý thuyết khoa học máy tính lẫn giá trị ứng dụng thực tiễn trong công nghiệp chế tạo. Về mặt lý thuyết, nghiên cứu hoàn thiện mô hình Mạng khí Nơ-ron tăng trưởng Tích lũy Bằng chứng với khả năng kiểm soát hiện tượng bùng nổ số lượng nút, ổn định vị trí trọng số và tự thích ứng với các dạng dữ liệu mới phát sinh theo thời gian \cite{fritzke1995growing}. Về mặt thực tiễn, luận văn cung cấp giải pháp phân loại lỗi ngoại quan tự động cho các dây chuyền lắp ráp màn hình điện thoại, giúp loại bỏ công đoạn gán nhãn thủ công tốn kém chi phí, giảm thiểu tỷ lệ bỏ sót khuyết tật và nâng cao năng suất vận hành toàn hệ thống v.v…

\section*{Kết luận Chương 1}
\addcontentsline{toc}{section}{Kết luận Chương 1}

Chương \ref{chap:c1} đã trình bày tổng quan về bài toán kiểm tra chất lượng bề mặt màn hình điện thoại trong môi trường sản xuất công nghiệp. Bề mặt màn hình điện thoại có cấu trúc phức tạp với nhiều lớp vật liệu khác nhau, khiến các dạng khuyết tật ngoại quan xuất hiện với các hình thái đa dạng. Các phương pháp kiểm tra truyền thống phụ thuộc nhiều vào việc gán nhãn thủ công hoặc các mô hình học máy giám sát cố định, gây ra khó khăn trong việc phát hiện và gom nhóm các dạng lỗi mới phát sinh. Nghiên cứu đã xác định việc kết hợp giữa phát hiện bất thường ngoại tuyến và phân cụm trực tuyến là hướng đi phù hợp để giải quyết vấn đề mất cân bằng dữ liệu và tự động phân loại khuyết tật.

Hệ thống đề xuất trong luận văn được xây dựng dựa trên kiến trúc hai pha nhằm tối ưu hóa tiến trình xử lý dòng dữ liệu ảnh. Pha phát hiện bất thường ngoại tuyến sử dụng bản đồ độ hỗn loạn thông tin không gian kết hợp với mô hình mã hóa tự động và thuật toán Mạng khí Nơ-ron tăng trưởng để nhanh chóng loại bỏ các sản phẩm đạt chất lượng. Pha phân cụm lỗi trực tuyến áp dụng thuật toán Mạng khí Nơ-ron tăng trưởng Tích lũy Bằng chứng cải tiến cùng mạng ResNet để tự động gom nhóm các dạng khuyết tật. Phương pháp đề xuất đã tích hợp bốn cơ chế kỹ thuật bao gồm (i) suy giảm tốc độ học theo độ tuổi của nút, (ii) chèn nút muộn dựa trên vùng đệm bằng chứng, (iii) kiểm soát sự phát triển hình thái dựa trên mật độ không gian, và (iv) gộp các vết lỗi dựa trên sự chồng lấp hộp giới hạn.

Các mục tiêu, nội dung nghiên cứu và kết quả dự kiến của luận văn đã được xác định rõ ràng nhằm đóng góp cả về mặt lý thuyết học máy lẫn tính ứng dụng thực tiễn trong công nghiệp. Luận văn hướng tới việc hoàn thiện mô hình học động trực tuyến có khả năng tự điều chỉnh cấu trúc đồ thị, nâng cao độ chính xác phân loại lỗi và giảm thiểu chi phí vận hành cho nhà máy sản xuất. Để làm cơ sở khoa học cho việc triển khai các phương pháp đề xuất trong các chương tiếp theo, Chương 2 sẽ trình bày chi tiết cơ sở lý thuyết nền tảng về mạng khí nơ-ron tăng trưởng, mô hình mã hóa tự động và lý thuyết Mạng khí Nơ-ron tăng trưởng truyền thống v.v…